\documentclass[../hippo.tex]{subfiles}
\begin{document}

\section{Technical Preliminaries}
\label{sec:background}

We collect here some technical background that will be used in presenting the
general HiPPO framework and in deriving specific HiPPO update rules.

\subsection{Orthogonal Polynomials}
\label{sec:op}
Orthogonal polynomials are a standard tool for working with function spaces~\cite{chihara,szego}.
Every measure $\mu$ induces a unique (up to a scalar) sequence of \emph{orthogonal polynomials} (OPs) $P_0(x), P_1(x), \dots$ satisfying $\mathrm{deg}(P_i) = i$ and
$ \langle P_i, P_j \rangle_{\mu} \defeq \int P_i(x) P_j(x) \d \mu(x) = 0 $ for all $i \neq j$.
This is the sequence found by orthogonalizing the monomial basis $\{x^i\}$ with Gram-Schmidt with respect to $\langle \cdot, \cdot \rangle_\mu$.
The fact that OPs form an orthogonal basis is useful because the optimal polynomial $g$ of degree $\deg(g) < N$
that approximates a function $f$ is then given by
\begin{align*}%
  \sum_{i=0}^{N-1} c_i P_i(x)/\|P_i\|_\mu^2 \qquad \text{where } c_i = \langle f, P_i \rangle_\mu = \int f(x) P_i(x) \d \mu(x).
\end{align*}
Classical OPs families comprise Jacobi (which include Legendre and Chebyshev
polynomials as special cases), Laguerre, and Hermite polynomials.
The Fourier basis can also be interpreted as OPs on the unit circle in the
complex plane.



\subsubsection{Properties of Legendre Polynomials}
\label{sec:legendre-properties}

\paragraph{Legendre polynomials}

Under the usual definition of the canonical Legendre polynomial $P_n$, they are
orthogonal with respect to the measure $\omega^{\mathrm{leg}} = \mathbf{1}_{[-1,1]}$:
\begin{equation}
  \label{eq:legendre}
  \frac{2n+1}{2} \int_{-1}^1 P_n(x)P_m(x) \dd x = \delta_{nm}
\end{equation}

Also, they satisfy
\begin{align*}%
  P_n(1) &= 1 \\
  P_n(-1) &= (-1)^{n}
  .
\end{align*}

\paragraph{Shifted and Scaled Legendre polynomials}
We will also consider scaling the Legendre polynomials to be orthogonal on the
interval $[0, t]$.
A change of variables on \eqref{eq:legendre} yields
\begin{align*}
  (2n+1) \int_0^t P_n\left( \frac{2x}{t}-1 \right) P_m \left( \frac{2x}{t} - 1 \right) \frac{1}{t} \dd x
  &=
  (2n+1) \int P_n\left( \frac{2x}{t}-1 \right) P_m \left( \frac{2x}{t} - 1 \right) \omega^{\mathrm{leg}}\left( \frac{2x}{t}-1 \right) \frac{1}{t} \dd x
  \\
  &=
  \frac{2n+1}{2} \int P_n(x) P_m(x) \omega^{\mathrm{leg}}(x) \dd x
  \\
  &= \delta_{nm}
  .
\end{align*}
Therefore, with respect to the measure $\omega_t = \mathbf{1}_{[0,t]}/t$ (which is a probability measure for all $t$),
the normalized orthogonal polynomials are
\begin{align*}
  (2n+1)^{1/2} P_n\left( \frac{2x}{t}-1 \right)
  .
\end{align*}

Similarly, the basis
\begin{align*}
  (2n+1)^{1/2} P_n\left( 2\frac{x-t}{\theta} + 1 \right)
\end{align*}
is orthonormal for the uniform measure $\frac{1}{\theta}\mathbb{I}_{[t-\theta, t]}$.

In general, the orthonormal basis for any uniform measure
consists of $(2n+1)^{\frac{1}{2}}$ times the corresponding linearly shifted version of $P_n$.

\paragraph{Derivatives of Legendre polynomials}

We note the following recurrence relations on Legendre polynomials (\citep[Chapter~12]{arfken2005mathematical}):
\begin{align*}
  (2n+1)P_n &= P_{n+1}' - P_{n-1}' \\
  P_{n+1}' &= (n+1)P_n + xP_n'
\end{align*}
The first equation yields
\begin{equation}%
  P_{n+1}' = (2n+1)P_n + (2n-3)P_{n-2} + \dots,
\end{equation}
where the sum stops at $P_0$ or $P_1$.

These equations directly imply
\begin{align}
  P_{n}' &= (2n-1)P_{n-1} + (2n-5)P_{n-3} + \dots
  \label{eq:legendre-d}
\end{align}
and
\begin{align}
  (x+1)P_n'(x) &= P_{n+1}' + P_n' - (n+1)P_n
  \nonumber \\
  &= nP_n + (2n-1)P_{n-1} + (2n-3)P_{n-2} + \dots
  .
  \label{eq:legendre-xd}
\end{align}
These will be used in the derivations of the HiPPO-LegT and HiPPO-LegS updates, respectively.

\subsubsection{Properties of Laguerre Polynomials}
\label{sec:laguerre-properties}

The standard Laguerre polynomials $L_n(x)$ are defined to be orthogonal with
respect to the weight function $e^{-x}$ supported on $[0, \infty)$, while the
generalized Laguerre polynomials (also called associated Laguerre polynomials)
$L_n^{(\alpha)}$ are defined to be orthogonal with respect to the weight function
$x^{\alpha} e^{-x}$ also supported on $[0, \infty)$:
\begin{equation}
  \label{eq:laguerre}
  \int_{0}^{\infty} x^\alpha e^{-x} L_n^{(\alpha)}(x) L_m^{(\alpha)}(x) \d x = \frac{(n+\alpha)!}{n!} \delta_{n, m}.
\end{equation}
Also, they satisfy
\begin{equation}%
  \label{eq:laguerre-endpoint}
  L_n^{(\alpha)}(0) = \binom{n + \alpha}{n} = \frac{\Gamma(n+\alpha+1)}{\Gamma(n+1)\Gamma(\alpha+1)}.
\end{equation}

The standard Laguerre polynomials correspond to the case of $\alpha = 0$ of
generalized Laguerre polynomials.

\paragraph{Derivatives of generalized Laguerre polynomials}

We note the following recurrence relations on generalized Laguerre polynomials (\citep[Chapter~13.2]{arfken2005mathematical}):
\begin{align*}
  \frac{\d}{\d x} L_n^{(\alpha)}(x) &= -L_{n-1}^{(\alpha+1)} (x) \\
  L_{n}^{(\alpha+1)}(x) &= \sum_{i=0}^{n} L_i^{(\alpha)}(x).
\end{align*}
These equations imply
\begin{equation*}
  \ddt L_n^{(\alpha)}(x) = - L_{0}^{(\alpha)}(x) - L_{1}^{(\alpha)} (x) - \dots - L_{n-1}^{(\alpha)}(x).
\end{equation*}

\subsubsection{Properties of Chebyshev polynomials}
\label{sec:chebyshev-properties}

Let $T_n$ be the classical Chebyshev polynomials (of the first kind), defined to
be orthogonal with respect to the weight function $(1-x^2)^{1/2}$ supported on
$[-1, 1]$,
and let $p_n$ be the normalized version of $T_n$ (i.e, with norm 1):
\begin{align*}
  \omega^{\mathrm{cheb}} &= (1-x^2)^{-1/2} \mathbb{I}_{(-1, 1)},
  \\
  p_n(x) &= \sqrt{\frac{2}{\pi}} T_n(x) \qquad \text{for } n \geq 1,
  \\
  p_0(x) &= \frac{1}{\sqrt\pi}.
\end{align*}
Note that $\omega^{\mathrm{cheb}}$ is not normalized (it integrates to $\pi$).


\paragraph{Derivatives of Chebyshev polynomials}
The chebyshev polynomials satisfy
\begin{align*}
  2T_n(x) = \frac{1}{n+1} \frac{d}{d x} T_{n+1}(x) - \frac{1}{n-1} \frac{d}{dx} T_{n-1}(x) \qquad n = 2, 3, \dots
  .
\end{align*}
By telescoping this series, we obtain
\begin{equation}
  \label{eq:chebyshev-derivative}
  \frac{1}{n} T_n' =
  \begin{cases}
    2(T_{n-1} + T_{n-3} + \dots + T_2) + T_0  & n \mbox{ odd}
    \\
    2(T_{n-1} + T_{n-3} + \dots + T_1) & n \mbox{ even}
  \end{cases}
  .
\end{equation}



\paragraph{Translated Chebyshev polynomials}
We will also consider shifting and scaling the Chebyshev polynomials to be orthogonal on the
interval $[t-\theta, t]$ for fixed length $\theta$.

The normalized (probability) measure is
\begin{align*}%
  \omega(t, x) &= \frac{2}{\theta\pi} \omega^{\mathrm{cheb}}\left( \frac{2(x-t)}{\theta} + 1 \right)
  = \frac{1}{\theta\pi} \left( \frac{x-t}{\theta} + 1 \right)^{-1/2} \left( -\frac{x-t}{\theta} \right)^{-1/2}
  \mathbb{I}_{(t-\theta, t)}
  .
\end{align*}
The orthonormal polynomial basis is
\begin{align*}%
  p_n(t, x) &= \sqrt\pi p_n\left( \frac{2(x-t)}{\theta} + 1 \right)
  .
\end{align*}
In terms of the original Chebyshev polynomials, these are
\begin{align*}%
  p_n(t, x) &= \sqrt{2} T_n\left( \frac{2(x-t)}{\theta} + 1 \right) \qquad \text{for } n \geq 1,
  \\
  p_0^{(t)} &= T_0\left( \frac{2(x-t)}{\theta} + 1 \right).
\end{align*}


\subsection{Leibniz Integral Rule}
As part of our standard strategy for deriving HiPPO update rules (\cref{sec:hippo-framework-details}),
we will differentiate through integrals with changing limits.
For example, we may wish to differentiate with respect to $t$ the expression
$\int f(t, x) \mu(t, x) \dd x = \int_0^t f(t, x) \frac{1}{t} \d x$
when analyzing the scaled Legendre (LegS) measure.

Differentiating through such integrals can be formalized by the Leibniz integral rule,
the basic version of which states that
\[
  \frac{\partial}{\partial t} \int_{\alpha(t)}^{\beta(t)} f(x, t) \dd x
  = \int_{\alpha(t)}^{\beta(t)} \frac{\partial}{\partial t} f(x, t) \dd x
  - \alpha'(t) f(\alpha(t), t) + \beta'(t) f(\beta(t), t).
\]

We elide over the formalisms in our derivations (\cref{sec:derivations}) and
instead use the following trick.
We replace integrand limits with an indicator function;
and using the Dirac delta function $\delta$ when differentiating (i.e., using the formalism
of distributional derivatives).
For example, the above formula can be derived succinctly with this trick:
\begin{align*}%
  \frac{\partial}{\partial t} \int_{\alpha(t)}^{\beta(t)} f(x, t) \dd x
  &= \frac{\partial}{\partial t} \int f(x, t) \mathbb{I}_{[\alpha(t), \beta(t)]}(x) \dd x
  \\
  &= \int \frac{\partial}{\partial t} f(x, t) \mathbb{I}_{[\alpha(t), \beta(t)]}(x) \dd x
  + \int f(x, t) \frac{\partial}{\partial t} \mathbb{I}_{[\alpha(t), \beta(t)]}(x) \dd x
  \\
  &= \int \frac{\partial}{\partial t} f(x, t) \mathbb{I}_{[\alpha(t), \beta(t)]}(x) \dd x
  + \int f(x, t) (\beta'(t) \delta_{\beta(t)}(x) - \alpha'(t) \delta_{\alpha(t)})(x) \dd x
  \\
  &= \int_{\alpha(t)}^{\beta(t)} \frac{\partial}{\partial t} f(x, t) \dd x
  - \alpha'(t) f(\alpha(t), t) + \beta'(t) f(\beta(t), t)
  .
\end{align*}


\subsection{ODE Discretization}
\label{sec:discretization-full}

In our framework, time series inputs will be modeled with a continuous function and then discretized.
Here we provide some background on ODE discretization methods,
including a new discretization that applies to a specific type of ODE that our new method encounters.

The general formulation of an ODE is $\frac{d}{dt} c(t) = f(t, c(t))$.
We will also focus on the linear time-invariant ODE of the form $\frac{d}{dt} c(t) = A c(t) + Bf(t)$
for some input function $f(t)$, as a special case.
The general methodology for discretizing the ODE, for step size $\Delta t$, is to
rewrite the ODE as
\begin{equation}
  \label{eq:ode-rhs}
  c(t + \Delta t) - c(t) = \int_t^{t+\Delta t} f(s, c(s)) \dd s,
\end{equation}
then approximate the RHS integral.

Many ODE discretization methods corresponds to different ways to
approximate the RHS integral:

\paragraph{Euler (aka forward Euler).}
To approximate the RHS of equation~\eqref{eq:ode-rhs}, keep the left endpoint
$\Delta t f(t, c(t))$.
For the linear ODE, we get:
\begin{equation*}
  c(t + \Delta t) = (I + \Delta t A) c(t)+ \Delta t B f(t).
\end{equation*}


\paragraph{Backward Euler.}
To approximate the RHS of equation~\eqref{eq:ode-rhs}, keep the right endpoint
$\Delta t f(t+\Delta t, c(t+\Delta t))$.
For the linear ODE, we get the linear equation and the update:
\begin{align*}
  c(t + \Delta t) - \Delta t A c(t + \Delta t) &= c(t)+ \Delta t B f(t) \\
  c(t + \Delta t) &= (I - \Delta t A)^{-1} c(t) + \Delta t (I - \Delta t A)^{-1} B f(t).
\end{align*}

\paragraph{Bilinear (aka Trapezoid rule, aka Tustin's method).}
To approximate the RHS of equation~\eqref{eq:ode-rhs}, average the endpoints
$\Delta t \frac{f(t, c(t))+f(t+\Delta t, c(t+\Delta t))}{2}$.
For the linear ODE, again we get a linear equation and the update:
\begin{align*}
  c(t + \Delta t) - \frac{\Delta t}{2} A c(t + \Delta t) &= (I + \Delta t/2 A) c(t) + \Delta t B f(t) \\
  c(t + \Delta t) &= (I - \Delta t/2 A)^{-1}(I + \Delta t/2 A) c(t) + \Delta t (I - \Delta t/2 A)^{-1} B f(t).
\end{align*}

\paragraph{Generalized Bilinear Transformation (GBT).}
This method~\citep{zhang2007performance} approximates the RHS of
equation~\eqref{eq:ode-rhs} by taking a weighted average of the endpoints
$\Delta t [(1 - \alpha) f(t, c(t)) + \alpha f(t+\Delta t, c(t+\Delta t))]$, for some parameter $\alpha \in [0, 1]$.
For the linear ODE, again we get a linear equation and the update:
\begin{align}
  c(t + \Delta t) - \Delta t \alpha A c(t + \Delta t) &= (I + \Delta t (1 - \alpha) A) c(t) + \Delta t B f(t) \nonumber \\
  c(t + \Delta t) &= (I - \Delta t \alpha A)^{-1}(I + \Delta t (1 - \alpha) A) c(t) + \Delta t (I - \Delta t \alpha A)^{-1} B f(t).
  \label{eq:gbt}
\end{align}
GBT generalizes the three methods mentioned above: forward Euler corresponds to
$\alpha = 0$, backward Euler to $\alpha = 1$, and bilinear to $\alpha = 1/2$.

We also note another method called \emph{Zero-order Hold}
(ZOH)~\citep{decarlo1989linear} that specializes to linear ODEs.
The RHS of equation~\eqref{eq:ode-rhs} is calculated in closed-form assuming
constant input $f$ between $t$ and $t + \Delta t$.
This yields the update
$c(t + \Delta t) = e^{\Delta t A} c(t) + \left( \int_{\tau=0}^{\Delta t} e^{\tau A} \d \tau \right) B f(t)$.
If $A$ is invertible, this can be simplified as
$c(t + \Delta t) = e^{\Delta t A} c(t) + A^{-1} (e^{\Delta t A} - I) B f(t)$.




\paragraph{HiPPO-LegS invariance to discretization step size.}
In the case of HiPPO-LegS, we have a linear ODE of the form
$\frac{d}{dt} c(t) = \frac{1}{t} A c(t) + \frac{1}{t} Bf(t)$.
Adapting the GBT discretization (which generalizes forward/backward Euler and
bilinear) to this linear ODE, we obtain:
\begin{align*}
  c(t + \Delta t) - \Delta t \alpha \frac{1}{t + \Delta t} A c(t + \Delta t) &= \left(I + \Delta t (1 - \alpha) \frac{1}{t} A\right) c(t) + \Delta t \frac{1}{t} B f(t) \\
  c(t + \Delta t) &= \left(I - \frac{\Delta t}{t + \Delta t}\alpha A\right)^{-1}\left(I + \frac{\Delta t}{t} (1 - \alpha) A\right) c(t)
               + \frac{\Delta t}{t} \left(I - \frac{\Delta t}{t + \Delta t} \alpha A\right)^{-1} B f(t).
\end{align*}
We highlight that this system is invariant to the discretization step size
$\Delta t$.
Indeed, if $c^{(k)} \defeq c(k \Delta t)$ and $f_k \defeq f(k \Delta t)$ then we have the recurrence
\begin{equation*}
  c^{(k+1)} = \left(I - \frac{1}{k+1}\alpha A\right)^{-1}\left(I + \frac{1}{k} (1 - \alpha) A\right) c^{(k)}
  + \frac{1}{k} \left(I - \frac{1}{k+1} \alpha A\right)^{-1} B f_k,
\end{equation*}
which does not depend on $\Delta t$.

\paragraph{Ablation: comparison between different discretization methods}
To understand the impact of approximation error in discretization, in
\cref{fig:discretization_error}, we show the absolute error for the HiPPO-LegS
updates in function approximation (\cref{sec:exp-hippo-ablations})
for different discretization methods: forward Euler, backward Euler,
and bilinear.
The bilinear method generally provide sufficiently accurate approximation.
We will use bilinear as the discretization method for the LegS updates
for the experiments.
\begin{figure}[ht]
  \centering
  \includegraphics[width=0.5\linewidth]{figures/legs_function_approx_error}
  \caption{Absolute error for different discretization methods. Forward and
    backward Euler are generally not very accurate, while bilinear
  yields more accurate approximation.}
  \label{fig:discretization_error}
\end{figure}

\end{document}

